\phantomsection
\section*{To the Reader}
\addcontentsline{toc}{section}{To the Reader}

Begin before Advent does. In the days before the season opens, read the first
part, The Coming, on what the Church is waiting for and how she learned to
wait. Then read the second part, Keeping Advent, and before the First Sunday
make a rule for the four weeks and settle when you will go to confession. The
third part, Through the Weeks, goes through the season in order: read each
week's chapter on the Sunday that opens it, the chapter on the Ember Days in
the third week, the chapter on the great antiphons on 17 December, and the
chapter on the Vigil on the evening before it. Keep the fourth part, Toward
Bethlehem, for the last days, so that on the holy night you know what you are
looking at. For every day of the season there is a reading in
Appendix~\ref{sec:daily}, and the prayers and hymns of the season, in Latin and
English, are in Appendix~\ref{sec:prayers}.

The Church prays through the day and the night in her Office, the Hours of
psalms, hymns, lessons and prayers, and the book that holds them is the
Breviary. Matins is the Office of the night, and in it the lessons from
Scripture and the Fathers are read. Lauds is the
praise of the dawn. Through the day come shorter Hours, among them None, the
ninth hour, in the middle of the afternoon. Vespers is the evening Office, and a Sunday or a great feast begins on the evening before, at its
First Vespers. Compline is the last Office of all, said before going to bed.
The book of the Mass is the Missal, with the prayers and readings of every day,
and the Mass's prayer of the day is called the Collect. Scripture in English is
quoted from the Douay-Rheims, the Catholic Bible translated from the Latin at
Rheims and Douay and revised by Bishop Challoner, so that its words follow the
Latin the Church reads at Mass and in the Office.
